% -----------------------------------------------------------------------------
% Yukawa 衰变与阈值定律：主文件保留目标式，计算统一进入证明区
% -----------------------------------------------------------------------------

同一个 Yukawa 顶角既能作为散射中的内部交换，也允许足够重的标量媒介粒子衰变为费米子--反费米子对。继续采用\eqnrefs{eq:yukawa-full-si-r10} 的无量纲耦合 $y$、物理四动量 $p^\mu=(E/c,\mathbf p)$ 和协变旋量归一化 $\bar uu=2mc$。

\begin{exbox}[标量 Yukawa 衰变]
设标量质量为 $m_Y$、费米子质量为 $m$。过程
\begin{equation*}
 \phi(P)\to f(p_1,s_1)+\bar f(p_2,s_2),
 \qquad P=p_1+p_2,
\end{equation*}
要求 $m_Y\ge2m$。树级约化振幅为
\begin{equation}
 \boxed{\mathcal M_S=-y\,\bar u(p_1)v(p_2).}
 \label{eq:yukawa-scalar-decay-M-si-r10}
\end{equation}
末态自旋求和得到
\begin{equation}
 \boxed{
 \sum_{s_1,s_2}|\mathcal M_S|^2
 =2y^2(m_Y^2-4m^2)c^2.}
 \label{eq:yukawa-scalar-decay-M2-si-r10}
\end{equation}
定义阈值速度参数
\begin{equation}
 \beta_S\equiv\sqrt{1-\frac{4m^2}{m_Y^2}},
 \qquad
 |\mathbf p_*|=\frac{m_Yc}{2}\beta_S,
 \label{eq:yukawa-decay-beta-si-r10}
\end{equation}
则静止系总衰变率为
\begin{equation}
 \boxed{
 \Gamma_{{\rm d},S}
 =\frac{y^2m_Y c^2}{8\pi\hbar}\,\beta_S^3.}
 \label{eq:yukawa-scalar-decay-width-si-r10}
\end{equation}
阈值附近的三个 $\beta_S$ 中，一个来自二体相空间，另外两个来自标量双线性型在阈值处的零点。


\medskip

\eqnrefs{eq:yukawa-scalar-decay-M2-si-r10} 表明标量衰变到费米子对的自旋求和振幅在阈值处已经含两个速度因子。二体相空间再提供一个 $\beta_{\rm rel}^{(S)}$；在母粒子静止系中显式完成能量 $\delta$ 函数的 Jacobi 化简后，三个因子相乘便给出\eqnrefs{eq:yukawa-scalar-decay-width-si-r10} 的 $\beta^3$ 阈值律。

由已给出的自旋求和，
\begin{equation*}
\sum|\mathcal M_S|^2
=2y^2(m_Y^2-4m^2)c^2
=2y^2m_Y^2c^2\bigl(\beta_{\rm rel}^{(S)}\bigr)^2,
\end{equation*}
其中
\begin{equation*}
\beta_{\rm rel}^{(S)}=\sqrt{1-\frac{4m^2}{m_Y^2}}.
\end{equation*}
剩下的一个 $\beta_{\rm rel}^{(S)}$ 来自二体相空间。母粒子静止系中
\begin{equation*}
P^\mu=(m_Yc,\mathbf0),
\qquad
p_1^\mu=(E_*/c,\mathbf p),
\qquad
p_2^\mu=(E_*/c,-\mathbf p),
\end{equation*}
能量 \(\delta\) 函数要求 $m_Yc^2=2E_*$。从
\begin{equation*}
E_*^2=m^2c^4+p^2c^2
\end{equation*}
解得
\begin{equation*}
p_*=\frac{m_Yc}{2}\beta_{\rm rel}^{(S)}.
\end{equation*}
使用全书二体衰变主式
\begin{equation*}
\Gamma
=\frac{1}{2m_Yc^2}\frac{1}{\hbar}
\int\dd\mathfrak P_2\,\sum|\mathcal M|^2,
\end{equation*}
以及已经由 $\delta$-雅可比行列式推出的
\begin{equation*}
\int\dd\mathfrak P_2
=\frac{p_*}{4\pi m_Yc}
=\frac{\beta_{\rm rel}^{(S)}}{8\pi},
\end{equation*}
逐项代入：
\begin{align*}
\Gamma_{{\rm d},S}
&=\frac{1}{2m_Yc^2}\frac1\hbar
\frac{\beta_{\rm rel}^{(S)}}{8\pi}
\left(2y^2m_Y^2c^2\bigl(\beta_{\rm rel}^{(S)}\bigr)^2\right)\\
&=\frac{y^2m_Yc^2}{8\pi\hbar}\bigl(\beta_{\rm rel}^{(S)}\bigr)^3.
\end{align*}
因此三次幂不是经验阈值指数：一个 $\beta_{\rm rel}^{(S)}$ 是可用相空间，另外两个来自标量双线性型在阈值处的零点。
\end{exbox}

\begin{exbox}[赝标量 Yukawa 阈值]
令赝标量质量为 $m_P$、无量纲耦合为 $y_P$，相互作用采用\eqnrefs{eq:pseudoscalar-yukawa-si-r10}。过程 $\phi_P\to f\bar f$ 的末态自旋求和为
\begin{equation}
 \boxed{
 \sum|\mathcal M_P|^2=2y_P^2m_P^2c^2.}
 \label{eq:yukawa-pseudo-decay-M2-si-r10}
\end{equation}
定义
\begin{equation}
 \beta_P\equiv\sqrt{1-\frac{4m^2}{m_P^2}},
 \qquad
 |\mathbf p_*|=\frac{m_Pc}{2}\beta_P,
 \label{eq:yukawa-pseudo-beta-si-r15}
\end{equation}
可得
\begin{equation}
 \boxed{
 \Gamma_{{\rm d},P}
 =\frac{y_P^2m_P c^2}{8\pi\hbar}\,\beta_P.}
 \label{eq:yukawa-pseudo-decay-width-si-r10}
\end{equation}
标量与赝标量共享同一二体相空间，因此阈值指数 $\beta_S^3$ 与 $\beta_P$ 的差异来自 Clifford 顶角结构，而不是相空间公式。


\medskip

赝标量顶角中的 $\gamma^5$ 改变了阈值处的自旋结构。\eqnrefs{eq:yukawa-pseudo-decay-M2-si-r10} 在阈值处不再额外消失，因此衰变宽度的唯一速度因子来自二体相空间；把同一相空间 Jacobi 因子代入后，就得到\eqnrefs{eq:yukawa-pseudo-decay-width-si-r10} 的线性阈值律。

令 $\mathcal O_P=\ii\gamma^5$。Dirac 伴随为
\begin{align*}
\overline{\mathcal O}_P
&=\gamma^0\mathcal O_P^\dagger\gamma^0
=\gamma^0(-\ii\gamma^5)\gamma^0\\
&=-\ii\gamma^0\gamma^5\gamma^0
=\ii\gamma^5,
\end{align*}
其中最后一步用了 $\{\gamma^5,\gamma^0\}=0$ 与 $(\gamma^0)^2=1$。自旋求和因而为
\begin{align*}
\sum|\mathcal M_P|^2
&=y_P^2\Tr[(\slashed p_1+mc)\ii\gamma^5
(\slashed p_2-mc)\ii\gamma^5]\\
&=-y_P^2\Tr[(\slashed p_1+mc)\gamma^5
(\slashed p_2-mc)\gamma^5].
\end{align*}
再逐项使用
\begin{equation*}
\gamma^5\slashed p_2\gamma^5=-\slashed p_2,
\qquad
\gamma^5(-mc)\gamma^5=-mc,
\end{equation*}
所以括号变成 $-\slashed p_2-mc$，外部负号抵消：
\begin{align*}
\sum|\mathcal M_P|^2
&=y_P^2\Tr[(\slashed p_1+mc)(\slashed p_2+mc)]\\
&=4y_P^2(p_1\!\cdot p_2+m^2c^2).
\end{align*}
由 $P=p_1+p_2$ 与 $P^2=m_P^2c^2$，
\begin{equation*}
p_1\!\cdot p_2=\frac{m_P^2-2m^2}{2}c^2,
\end{equation*}
于是
\begin{equation*}
\sum|\mathcal M_P|^2=2y_P^2m_P^2c^2.
\end{equation*}
它在阈值 $m_P\to2m$ 时并不再额外消失，因此唯一的速度因子来自二体相空间：
\begin{equation*}
\int\dd\mathfrak P_2=\frac{\beta_{\rm rel}^{(P)}}{8\pi},
\qquad
\beta_{\rm rel}^{(P)}=\sqrt{1-\frac{4m^2}{m_P^2}}.
\end{equation*}
代入衰变率主式，
\begin{equation*}
\Gamma_{{\rm d},P}
=\frac{y_P^2m_Pc^2}{8\pi\hbar}\beta_{\rm rel}^{(P)}.
\end{equation*}
因此标量的 $[\beta_{\rm rel}^{(S)}]^3$ 与赝标量的 $\beta_{\rm rel}^{(P)}$ 阈值幂次差别完全来自 Clifford 结构对阈值矩阵元的不同约束，而非相空间公式不同。
\end{exbox}

\begin{figure}[htbp]
 \centering
 \includegraphics[width=.72\linewidth]{yukawa_decay_threshold.png}
 \caption{Yukawa 媒介粒子的阈值线型。纵轴为 $8\pi\hbar\Gamma_{\rm d}/(y_X^2m_Xc^2)$，其中 $X=S,P$；标量耦合给出 $\beta_S^3$，赝标量耦合给出 $\beta_P$。}
 \label{fig:yukawa-decay-threshold-r6}
\end{figure}

开放衰变通道 $\phi\to f\bar f$ 同时使标量二点函数产生自能虚部。把一般共振展开~\eqnrefs{eq:resonance-pole-SI-r12} 专门化到这一通道时，能量宽度为
\begin{equation*}
 \Gamma_E=\hbar\Gamma_{\rm d},
\end{equation*}
其中 $\Gamma_{\rm d}$ 的 SI 单位为 $\mathrm{s^{-1}}$。在\eqnrefs{eq:resonance-pole-SI-r12} 的极点隔离与背景抑制条件同时成立时，极点项的虚能量位移为 $\hbar\Gamma_{\rm d}/2$。相空间计算得到的衰变率与光学定理给出的在壳自能虚部，是同一开放通道的两种表示。

